% !TeX root = ../navier-stokes-manuscript.tex
% ============================================================
% 2.2 Finite-Time Response Escalation
% ============================================================

\subsection{What finite time implies after one response escapes}\label{sub:FiniteTimeResponseEscalation}

The equation supplies the temporal rows; finite time supplies the implication
between two adjacent rows.  If one response has no finite roof while the next is
integrable, the fundamental theorem of calculus would nevertheless give the
first a terminal limit.  The following argument records that contradiction.  It
propagates escape through rows already known to exist and creates none of them.

Use one fixed Banach space \(B\) throughout and write the Eulerian rows as
\[
  V_n(t):=\partial_t^n u(\cdot,t),
\]
Assume \(V_n:[0,T_*)\to B\) is locally absolutely continuous and has Bochner
derivative \(V_{n+1}\).  For \(0\le s<t<T<T_*\),
\[
  V_n(t)-V_n(s)
  =
  \int_s^t V_{n+1}(\tau)\,d\tau
\]
and therefore
\[
  \lVert V_n(t)-V_n(s)\rVert_B
  \le
  \int_s^t \lVert V_{n+1}(\tau)\rVert_B\,d\tau.
\]
Suppose
\[
  \sup_{0<t<T_*}\lVert V_n(t)\rVert_B
  =
  \infty,
\]
Since \(\lVert V_n(0)\rVert_B\) is finite, the integral on the right cannot
remain finite; hence
\[
  \int_0^{T_*}\lVert V_{n+1}(\tau)\rVert_B\,d\tau
  =
  \infty.
\]
The interval itself is finite, so an \(L^1\) divergence forces
\[
  \sup_{0<t<T_*}\lVert V_{n+1}(t)\rVert_B
  =
  \infty.
\]
Iteration is licensed only when every successive pair has this same fixed
space, local absolute continuity, and Bochner-derivative identity.  Under those
hypotheses, finite time turns escape of one temporal response into escape of the
next.
